\section{DIAGRAMAS DE FEYNMAN (FD): GENERALIDADES}
Representación gráfica en el \textbf{espacio de momentos}. \\
$\bullet$ \text{Ejes:} Eje temporal de izquierda (IN) a derecha (OUT).\\
$\bullet$ Conservación $P^\mu$ en cada vértice y globalmente.\\
$\bullet$ Líneas Internas: Virtuales (off-shell: $p^2 \neq m^2$).\\
$\bullet$ Líneas externas (patas): Reales (on-shell: $p_i^2 = m_i ^2$).\\
$\bullet$ Amplitud Invariante $(F \text{ o } \mathcal{M})$:
$$T_{fi} = - i (2\pi)^4 \delta^{(4)}(P_{\text{tot}}) F \qquad P_{\text{tot}}= P_{\text{in}} - P_{\text{out}}$$
$$ -iF = (\text{Vértices}) \times (\text{Propagadores}) \times (\text{Patas Ext.}) $$


\section{GUÍA DE TRAZADO LÍNEAS}
\begin{center}
\begin{tabular}{cll}
\textbf{Trazo} & \textbf{Estilo Visual} & \textbf{Partícula} \\
\hline
\tikz[baseline=-3pt]{\draw[fermion] (0,0) -- (1,0);} & Línea cont. + Flecha & \textbf{Fermión} ($e^-, \mu, q,  \nu$) \\
\tikz[baseline=-3pt]{\draw[antifermion] (0,0) -- (1,0);} & Flecha opuesta & \textbf{Antifermión} ($e^+, \bar{q}$) \\
\tikz[baseline=-3pt]{\draw[photon] (0,0) -- (1,0);} & Ondulada (Snake) & \textbf{Fotón} ($\gamma, W, Z$) \\
\tikz[baseline=-3pt]{\draw[scalar] (0,0) -- (1,0);} & Discontinua (Dashed) & \textbf{Escalar} ($\pi, K,$\textit{Higgs}) \\
\tikz[baseline=-3pt]{\draw[gluon] (0,0) -- (1,0);} & Bucle (Coil) & \textbf{Gluón} ($g$) \\
\end{tabular}\\
$\bullet$ \textbf{Flechas}: Indican flujo de carga fermiónica (NO siempre coinciden con el momento).
\end{center}
\section{CANALES DE MANDELSTAM ($1+2 \to 3+4$)}
Amplitud total $F = F_s + F_t + F_u$. \quad\quad $s+t+u = m_A^2+ m_B^2 + m_C^2 + m_D^2$
\begin{center}
\begin{tikzpicture}[scale=0.5, baseline=(current bounding box.center)]
\node at (1.5,2.5) {\textbf{Canal s}: Aniquilación};
\draw[fermion] (0,2) node[left]{$p_1$} -- (1,1);
\draw[fermion] (0,0) node[left]{$p_2$} -- (1,1);
\draw[photon] (1,1) -- (2.5,1) node[midway, above]{$s$};
\draw[fermion] (2.5,1) -- (3.5,2) node[right]{$p_3$};
\draw[fermion] (2.5,1) -- (3.5,0) node[right]{$p_4$};
\node at (1.5, -1) {\scalebox{0.8}{$s = (p_1+p_2)^2 =(p_3+p_4)^2$}};
\end{tikzpicture}
\hspace{5pt}
\begin{tikzpicture}[scale=0.5, baseline=(current bounding box.center)]
\node at (1.5,2.5) {\textbf{Canal t}: Dispersión};
\draw[fermion] (0,2) node[left]{$p_1$} -- (1,2);
\draw[fermion] (1,2) -- (3,2) node[right]{$p_3$};
\draw[photon] (1,2) -- (1,0) node[midway, right]{$t$};
\draw[fermion] (0,0) node[left]{$p_2$} -- (1,0);
\draw[fermion] (1,0) -- (3,0) node[right]{$p_4$};
\node at (1.5, -1) {\scalebox{0.8}{$t = (p_1-p_3)^2=(p_2-p_4)^2$}};
\end{tikzpicture}
\hspace{5pt}
\begin{tikzpicture}[scale=0.5, baseline=(current bounding box.center)]
\node at (1.5,2.5) {\textbf{Canal u}: Intercambio};
\draw[fermion] (0,2) node[left]{$p_1$} -- (1,2);
\draw[fermion] (1,0) -- (3,2) node[right]{$p_3$};
\draw[photon] (1,2) -- (1,0) node[midway, right]{$u$};
\draw[fermion] (0,0) node[left]{$p_2$} -- (1,0);
\draw[fermion] (1,2) -- (3,0) node[right]{$p_4$};
\node at (1.5, -1) {\scalebox{0.8}{$u = (p_1-p_4)^2=(p_2-p_3)^2$}};
\end{tikzpicture}
\end{center}
CM, $m \to 0$:\;\; $s = 4E^2 $\qquad\;\;\qquad\qquad$t= -\frac{s}{2}(1 - \cos\theta)$\qquad\qquad $u= -\frac{s}{2}(1 + \cos\theta)$ 

\section{REGLAS: PROPAGADORES (Internos)}
Inverso de la función de Green de 2 puntos en el espacio de momentos.\\
$\mathcal{FT}$ al espacio de momentos: $i\partial^\mu \rightarrow p^\mu \quad  \square^2 = \partial_\mu \partial^\mu \rightarrow -p^2\quad \frac{1}{\square^2 + m^2} \rightarrow\frac{1}{-p^2 + m^2} $\\
\textbf{1. Escalar} (spin 0, $m$): $\Delta(p) = \frac{i}{p^2 - m^2}$\\
\textbf{2. Fermión} (spin $1/2$, $m$): $S_F(p) = \frac{i(\cancel{p} + m)}{p^2 - m^2}$\\
\textbf{3. Fotón} (spin 1, $m=0$, Feynman): $D_{\mu\nu}(p) = \frac{-i g_{\mu\nu}}{p^2} = \frac{-i g_{\nu \mu}}{k^2}$

\section{REGLAS: VÉRTICES DE INTERACCIÓN -- $V^\mu = FR \equiv$ Feymann Rule}
Multiplica $i$ por $\mathcal{L}_{\text{int}}$ (solamente la parte relevante de $\mathcal{L}_{\text{int}}$ en el espacio de momentos: no normalización y no campos externos: e.g. $A_\mu$, $e^{i(p_{\text{f}}- p_{\text{i}})x}, \dots$).\\

\begin{minipage}{0.49\linewidth}
\textbf{1. Escalar:} Bosón cargado $\phi$ - Fotón $A_\mu$ 
\begin{center}
$\mathcal{L}_{\text{int}} \propto J^\mu A_\mu$\vspace{3pt}
\begin{tikzpicture}[scale=0.6, baseline=(current bounding box.center)]
\draw[scalar, ->] (-1,0.8) -- (0,0) node[near end, above left]{$p_{\text{in}}\quad$};
\draw[scalar, ->] (0,0) -- (1,0.8) node[near start, above right]{$\quad p_{\text{out}}$};
\draw[photon] (0,0) -- (0,-1) node[below]{$\mu$};
\node[vertex] at (0,0) {};
\node at (3,-0.2) {$V^\mu = -i Q e (p_{\text{in}} + p_{\text{out}})^\mu$};
\end{tikzpicture}
\end{center}
\end{minipage}
\begin{minipage}{0.49 \linewidth}
\textbf{2. Fermiónica:} Fermión $\psi$ - Fotón $A_\mu$ 
\begin{center}
$\mathcal{L}_{\text{int}} = -e Q \bar{\psi} \gamma^\mu \psi A_\mu$\vspace{3pt}
\begin{tikzpicture}[scale=0.6, baseline=(current bounding box.center)]
\draw[fermion] (-1,0.8) -- (0,0);
\draw[fermion] (0,0) -- (1,0.8);
\draw[photon] (0,0) -- (0,-1) node[below]{$\mu$};
\node[vertex] at (0,0) {};
\node at (2.5,0) {$V^\mu = -i Q e \gamma^\mu$};
\end{tikzpicture}
\end{center}    
\end{minipage}



\section{TEOREMAS DE TRAZAS (TRUCO DE CASIMIR)}
$\Gamma$ cualquier producto de matrices de Dirac. $\bar{\Gamma} = \gamma^0 \Gamma^\dagger\gamma^0$\\
$\bullet$ \textbf{Relación}: $[\bar{u}(p',s')\, \Gamma\, u(p,s)]^\dagger = \bar{u}(p,s)\, \bar{\Gamma}\, u(p',s').$ 
\\ $\bullet$ If $\Gamma=\gamma^\mu:$ $[\bar{u}(p',s')\, \gamma^\mu\, u(p,s)]^\dagger = \bar{u}(p,s)\, \gamma^\mu\, u(p',s').$\\
$\bullet$ \textbf{Tip}: $\sum_{s_1, s_3}| \bar{u}_3 \gamma^\mu u_1 |^2 = \sum_{s_1, s_3}(\bar{u}_3 \gamma^\mu u_1)(\bar{u}_1 \gamma^\nu u_3)= \text{Tr}[(\cancel{p}_3 + m_3)\gamma^\mu (\cancel{p}_1 + m_1)\gamma^\nu]$.





\section{CÁLCULO DE SECCIÓN EFICAZ ($\sigma$)}
Relación entre Amplitud Invariante $F$ (o $\mathcal{M}$) y observables.\\
$\bullet$ \textbf{Sección Eficaz Diferencial} ($A+B \to C+D$ en CM):
$$ d\sigma = \frac{1}{2E_{CM} |\vec{v}_A - \vec{v}_B|} |\bar{F}|^2 d\Phi_2= \frac{1}{2\lambda^{1/2} (s, m_A^2, m_B^2)} |\vec{F}|^2 \, dLips(s, p_C, p_D) $$
$$\lambda^{1/2} (s, m_A^2, m_B^2) = 2 \sqrt{(p_A \cdot p_B)^2 - m_A^2 m_B^2}$$
$$ d{Lips}(s, p_C, p_D)  = (2\pi)^4 \delta^4(p_A+p_B- p_C -p_D) \frac{d^3\vec{p}_C}{(2\pi)^3 2E_C} \frac{d^3\vec{p}_D}{(2\pi)^3 2E_D}$$

$$\delta^4 (p_1 +p_2-p_3 -p_4) = \delta (E_1+E_2-E_3-E_4) \delta^3(\vec{p_1}+\vec{p_2}-\vec{p_3}-\vec{p_3})$$

$$ \textbf{Promedio de spin} \text{ (si no son escalares):}\quad|\bar{F}|^2 = \frac{1}{(2S_A+1)(2S_B+1)} \sum_{\text{spins}} |F|^2 $$
$$\text{Si la partícula $i$ es un \textbf{fotón}, no se emplea $(2S_i +1),$ sino que se divide entre 2}$$
$$ \frac{d\sigma}{d\Omega_{CM}} = \frac{1}{64\pi^2 s} \frac{|\vec{p}_f|}{|\vec{p}_i|} |\bar{F}|^2 \qquad \frac{d\sigma}{dt}= \frac{1}{16\pi} \frac{1}{\lambda(s, m_1^2, m_2^2)} |\bar{F}|^2 $$





% \section{EJEMPLO DISPERSIÓN $e^- \mu^- \to e^- \mu^-$ (QED)}
% \textbf{Diagrama (Canal t)}: Único diagrama (partículas distinguibles).
% \begin{center}
% \begin{tikzpicture}[scale=0.6]
% \draw[fermion] (-1,1) node[left]{$e^-(p_1)$} -- (0,0.5);
% \draw[fermion] (0,0.5) -- (1,1) node[right]{$e^-(p_3)$};
% \draw[photon] (0,0.5) -- (0,-0.5) node[midway, right]{$q$};
% \draw[fermion] (-1,-1) node[left]{$\mu^-(p_2)$} -- (0,-0.5);
% \draw[fermion] (0,-0.5) -- (1,-1) node[right]{$\mu^-(p_4)$};
% \node at (0.0, 0.9) {$\mu$};
% \node at (0.0, -0.9) {$\nu$};
% \end{tikzpicture}
% \end{center}
% $$ -iF = \underbrace{\{\bar{u}(p_3) (+ie\gamma^\mu) u(p_1)\}}_{\text{Corriente } e^-} \underbrace{\left(\frac{-ig_{\mu\nu}}{q^2}\right)}_{\text{Prop. Fotón}} \underbrace{\{\bar{u}(p_4) (+ie\gamma^\nu) u(p_2)\}}_{\text{Corriente } \mu^-} $$
% $$ F = -\frac{e^2}{t} [\bar{u}_3 \gamma^\mu u_1] [\bar{u}_4 \gamma_\mu u_2] $$
% $$ |\bar{F}|^2 = \frac{e^4}{4t^2} \text{Tr}[(\cancel{p}_3+m_e)\gamma^\mu(\cancel{p}_1+m_e)\gamma^\nu] \text{Tr}[(\cancel{p}_4+m_\mu)\gamma_\mu(\cancel{p}_2+m_\mu)\gamma_\nu] $$
% \textbf{Resultado (Límite Ultrarelativista $m \to 0$)}:
% Usando Mandelstam ($s, t, u$):
% $$ |\bar{F}|^2 \approx \frac{2e^4}{t^2} (s^2 + u^2) \qquad \frac{d\sigma}{d\Omega} = \frac{\alpha^2}{2s} \frac{s^2+u^2}{t^2} $$

\section{PROCESOS QED (DIAGRAMAS)}

\textbf{1. Aniquilación $e^- e^+ \to \mu^- \mu^+$}:
Solo canal s (carga total 0, muones distinguibles).
\begin{center}
\begin{tikzpicture}[scale=0.5]
\draw[fermion] (-1,1) node[left]{$e^-$} -- (0,0);
\draw[antifermion] (-1,-1) node[left]{$e^+$} -- (0,0);
\draw[photon] (0,0) -- (1.5,0) node[midway, above]{$s$};
\draw[fermion] (1.5,0) -- (2.5,1) node[right]{$\mu^-$};
\draw[antifermion] (1.5,0) -- (2.5,-1) node[right]{$\mu^+$};
\end{tikzpicture}
\end{center}

\textbf{2. Dispersión Møller $e^- e^- \to e^- e^-$}:
Partículas iguales. Interferencia Canal t + Canal u.
\begin{center}
\begin{tikzpicture}[scale=0.45]
\node at (0,1.2) {\textbf{Canal t}};
\draw[fermion] (-1,0.8) -- (0,0.4); \draw[fermion] (0,0.4) -- (1,0.8);
\draw[photon] (0,0.4) -- (0,-0.4);
\draw[fermion] (-1,-0.8) -- (0,-0.4); \draw[fermion] (0,-0.4) -- (1,-0.8);
\end{tikzpicture}
\hspace{0.2cm}
\begin{tikzpicture}[scale=0.45]
\node at (0,1.2) {\textbf{Canal u}};
\draw[fermion] (-1,0.8) -- (0,0.4); 
\draw[fermion] (0,-0.4) -- (1,0.8); % Cruzado
\draw[photon] (0,0.4) -- (0,-0.4);
\draw[fermion] (-1,-0.8) -- (0,-0.4); 
\draw[fermion] (0,0.4) -- (1,-0.8); % Cruzado
\end{tikzpicture}
\end{center}
$$ F = F_t - F_u \quad (\text{Signo relativo - por Fermi}) $$

\textbf{3. Dispersión Bhabha $e^- e^+ \to e^- e^+$}:
Canal s (aniquilación) + Canal t (scattering).
\begin{center}
\begin{tikzpicture}[scale=0.45]
\node at (0.8,1.2) {\textbf{Canal s}};
\draw[fermion] (-1,0.8) -- (0,0);
\draw[antifermion] (-1,-0.8) -- (0,0);
\draw[photon] (0,0) -- (1.5,0);
\draw[fermion] (1.5,0) -- (2.5,0.8);
\draw[antifermion] (1.5,0) -- (2.5,-0.8);
\end{tikzpicture}
\hspace{0.2cm}
\begin{tikzpicture}[scale=0.45]
\node at (0,1.2) {\textbf{Canal t}};
\draw[fermion] (-1,0.8) -- (0,0.4); \draw[fermion] (0,0.4) -- (1,0.8);
\draw[photon] (0,0.4) -- (0,-0.4);
\draw[antifermion] (-1,-0.8) -- (0,-0.4); \draw[antifermion] (0,-0.4) -- (1,-0.8);
\end{tikzpicture}
\end{center}

\textbf{4. Compton $e^- \gamma \to e^- \gamma$}:
Canal s + Canal u (Diagramas con propagador fermiónico).



